\section{引言}
脂肪酶催化酯化为生物质来源脂肪酸转化为润滑剂提供了反应路线。三羟甲基丙烷（TMP）经逐级酯化生成单酯（M）、二酯（D）和三酯（T），产物组成影响润滑剂的性能，因而反应器设计需要同时描述转化速率与酯类分布 \cite{akerman2011,mao2024}。近年来，Montiel 等研究了无溶剂条件下支链酸与支链醇的酶催化生物润滑剂合成及催化剂重复使用 \cite{montiel2024}；Mao 等将油酸裂解所得壬酸用于 TMP 酯合成，并评价了脂肪酶批次稳定性和产物润滑性能 \cite{mao2024}。可再生资源制备 TMP 的研究进一步拓展了全生物基原料路线 \cite{sayed2025}。合成路线与催化剂性能的进展使动力学模型的预测能力成为过程放大与条件选择的重要依据。

固定化酶反应的表观动力学还受到载体传质、底物吸附与酶稳定性的共同影响。Sengupta 等将竞争吸附和多孔树脂表征结合，研究了固定化脂肪酶催化酯化的吸附与孔扩散条件 \cite{sengupta2026}。Ribeiro 等在壳聚糖--黄原胶双层纳米胶囊固定化脂肪酶中考察了热稳定性，并用一阶模型描述主要热失活过程 \cite{ribeiro2025}。固定化脂肪酶的吸附与热失活研究分别说明了传质条件和活性衰减在特定固定化酶体系中的作用，但不能直接确定 TMP 连续酯化中的伪持续性概率或失活常数。TMP 动力学分析需要利用组分轨迹区分组成路径偏差与反应时间尺度偏差。

三个连续酯化步骤可由共享分母的 Ping Pong bi bi 方程描述 \cite{cleland1963,cornishbowden2012,bornadel2013}。完整速率方程包含多个相互关联的参数，有限浓度--时间观测往往不足以约束全部参数 \cite{alhaque2012}。Nishimura 等提出了两步半机理估计方法，本文称为初始两步法 \cite{nishimura2026}。指纹步骤以反应度 $c$ 为自变量，估计第二和第三次酯化的选择性 $S_M$、$S_D$；动态步骤固定选择性，再估计集总速率参数。公共速率因子在反应度域中消去，是指纹方法分离组成路径与时间尺度的基础 \cite{mitchell2022a,mitchell2022b,mitchell2023}。Nishimura 等利用四项研究的公开曲线，以残差平方和和贝叶斯信息准则（BIC）比较了简化与完整动态模型。

初始两步法的选择性约束仍存在两个需要检验的问题。不可逆指纹模型未包含伪持续性，即新生成的酯在释放前再次酰化的过程。Nishimura 等已指出，Bornadel 数据的传质限制可能使指纹参数仅代表表观选择性 \cite{nishimura2026,mitchell2023}。固定存在模型偏差的指纹估计，会限制动态步骤修正组分轨迹的能力。初始动态模型采用批次独立的速率常数描述重复使用过程，但未显式估计批内失活。此外，Nishimura 等的模型比较主要基于样本内拟合，BIC 参数计数未包括第一步从相同数据估计的两个选择性。组成路径修正能否改善留出观测预测，以及参数增加是否造成过拟合，需要独立于全数据拟合的检验。

近期酶反应网络与科学机器学习研究为参数估计和验证提供了相关思路。van Sluijs 等将最优实验设计与流动化学结合，通过迭代构建训练数据提高酶反应网络的动力学信息量 \cite{vansluijs2024}。Giampiccolo 等将部分已知动力学嵌入混合神经常微分方程，同时考察机理参数估计与可辨识性 \cite{giampiccolo2024}。de Rooij 等使用生理约束正则化改善通用微分方程在稀疏数据下的训练与预测 \cite{derooij2025}。酶反应网络实验设计与混合动力学研究表明，数据覆盖范围、机理约束和参数辨识需要联合评估，较低训练误差本身不能替代预测验证。表格基础模型 TabPFN 和 TabICL 也扩展了小样本数据建模的方法选择 \cite{hollmann2025,qu2025tabicl}，但通用表格预测能力并不直接等同于受质量守恒约束的反应外推能力。

本文研究指纹选择性约束与机理缺项对 TMP 酯化预测的影响，并提出带指纹信息先验的伪持续性动力学模型（PPK-FIP）。PPK-FIP 保留反应度域的选择性估计，将固定选择性改为经过校准的先验，并在动力学方程中显式加入伪持续性和一阶失活。精确时间重参数化用于降低多起点估计与交叉验证的计算成本。表~\ref{tab:roadmap} 概括了初始模型评估、组分残差诊断、样本外预测、组件消融、参数恢复及混合建模的检验设计。所有方法在相同公开数据上比较，指纹估计在每个验证折内重新计算；跨条件迁移和已知真值模拟用于限定机理参数的解释范围。

\begin{table}[!htbp]
\centering\small
\caption{TMP 酯化动力学模型的六项检验及对应证据。图~\ref{fig:workflow}、\ref{fig:data} 和 \ref{fig:bench} 给出估计流程、数据与计算性能。}\label{tab:roadmap}
\begin{tabular}{@{}c p{5.1cm} p{4.4cm} p{4.4cm}@{}}
\toprule
步骤 & 检验目标 & 分析方法 & 对应图表 \\
\midrule
1 & 初始两步模型的拟合一致性、参数可辨识性与样本内评价局限 & 公开数据核对、两步参数重新估计、多起点搜索与剖面目标函数分析 & 表~\ref{tab:data}--\ref{tab:reprodyn}；图~\ref{fig:repro}、\ref{fig:landscape} \\
2 & 初始两步模型的组分浓度残差与伪持续性、批内失活的关系 & 全数据拟合、组成路径比较与组分残差诊断 & 表~\ref{tab:insample}、\ref{tab:params}；图~\ref{fig:insample}--\ref{fig:resid} \\
3 & PPK-FIP 对留出时间点及后期反应轨迹的预测能力 & 留一时间点交叉验证与前半段训练的向前外推 & 表~\ref{tab:loto}、\ref{tab:fcst}；图~\ref{fig:loto}--\ref{fig:fcst} \\
4 & 伪持续性、失活、指纹先验及模型平均对预测误差的贡献 & 组件消融、先验宽度敏感性与跨条件参数池化 & 表~\ref{tab:ablation}、\ref{tab:pooling}；图~\ref{fig:ablation}、\ref{fig:pooling} \\
5 & 伪持续性组成路径的跨条件适用性与选择性估计偏差 & 留出条件的组成路径预测与蒙特卡洛参数恢复 & 表~\ref{tab:transfer}、\ref{tab:sim}；图~\ref{fig:transfer}、\ref{fig:sim} \\
6 & 机理约束与高斯过程残差校正对内插和外推的作用 & 数据驱动基线对比与机理--高斯过程混合建模 & 表~\ref{tab:lotobb}--\ref{tab:hybrid}；图~\ref{fig:blackbox}、\ref{fig:hybrid} \\
\bottomrule
\end{tabular}
\end{table}
